% GENERATED FILE. Edit canonical lessons, then run npm run generate:books.
% canonical-source-hash: fnv1a64:346d923518a77529
% canonical-lessons: FR-C398-reseau, FR-C398-connexion, FR-C398-naviguer, FR-C398-lien, FR-C398-publier

\chapter{Sur internet}
\label{ch:fr-sur-internet}
\hlchaptermodality{398}

\begin{chapteropening}
\textbf{By the end of this chapter:} \emph{I can talk about being online: the network, the connection, browsing, a link and posting something.}

Complete the last lesson of chapter 398: I can talk about being online: the network, the connection, browsing, a link and posting something.
\end{chapteropening}

\section[le réseau]{le réseau — the network}

\label{lesson:FR-C398-reseau}



\begin{quote}
Before the new one: say the French for to paste, then the French for to save.
\end{quote}

\subsection*{You'll want to know: le réseau}
\textbf{le réseau} — \textquotedblleft{}the network\textquotedblright{}.

\begin{quote}
Il n'y a pas de réseau en montagne.
\end{quote}

\begin{quote}
\emph{There is no signal in the mountains.}
\end{quote}

\emph{Les réseaux sociaux} — social media.

\begin{cousinweb}
\emph{réseau} is usually traced to Old French \emph{reseuil}, a small form of \emph{rets} \textquotedblleft{}net\textquotedblright{} (EtymDB has \emph{rets} from Old French \emph{reiz}), which goes back to Latin \emph{rēte} \textquotedblleft{}a net\textquotedblright{} (Lewis \& Short).
\end{cousinweb}

\begin{grammarlens}[title={What you've built}]
One more word for telling, explaining and describing at B1.
\end{grammarlens}

\subsection*{Your turn}
\begin{itemize}
\raggedright
  \item \emph{Say it:} \emph{le réseau}
  \item \emph{Say it:} \emph{le réseau}, once more
  \item \emph{Say it:} a sentence of your own with \emph{le réseau}
  \item \emph{Recall:} say the French for to paste, then the French for to save, then say \emph{le réseau} again
\end{itemize}

\begin{tcolorbox}[breakable,colback=teal!4,colframe=teal!35!black,title={Before you move on}]
What does le réseau mean? (\textquotedblleft{}The network\textquotedblright{}.) Say it once more.
\end{tcolorbox}
\section[la connexion]{la connexion — the connection}

\label{lesson:FR-C398-connexion}



\begin{quote}
Before the new one: say the French for to save, then the French for the network.
\end{quote}

\subsection*{You'll want to know: la connexion}
\textbf{la connexion} — \textquotedblleft{}the connection\textquotedblright{}.

\begin{quote}
La connexion est très lente ce soir.
\end{quote}

\begin{quote}
\emph{The connection is very slow tonight.}
\end{quote}

Spelled with \emph{x}.

\begin{cousinweb}
\emph{connexion} is a learned borrowing from Latin \emph{cōnexiō} \textquotedblleft{}a joining, a connection\textquotedblright{} (EtymDB), from \emph{cōnectere} \textquotedblleft{}to bind together\textquotedblright{}, built on \emph{nectere} \textquotedblleft{}to tie\textquotedblright{} (Lewis \& Short).
\end{cousinweb}

\begin{grammarlens}[title={What you've built}]
One more word for telling, explaining and describing at B1.
\end{grammarlens}

\subsection*{Your turn}
\begin{itemize}
\raggedright
  \item \emph{Say it:} \emph{la connexion}
  \item \emph{Say it:} \emph{la connexion}, once more
  \item \emph{Say it:} a sentence of your own with \emph{la connexion}
  \item \emph{Recall:} say the French for to save, then the French for the network, then say \emph{la connexion} again
\end{itemize}

\begin{tcolorbox}[breakable,colback=teal!4,colframe=teal!35!black,title={Before you move on}]
What does la connexion mean? (\textquotedblleft{}The connection\textquotedblright{}.) Say it once more.
\end{tcolorbox}
\section[naviguer]{naviguer — to browse; to sail}

\label{lesson:FR-C398-naviguer}



\begin{quote}
Before the new one: say the French for the network, then the French for the connection.
\end{quote}

\subsection*{You'll want to know: naviguer}
\textbf{naviguer} — \textquotedblleft{}to browse; to sail\textquotedblright{}.

\begin{quote}
Il navigue sur internet pendant des heures.
\end{quote}

\begin{quote}
\emph{He browses the internet for hours.}
\end{quote}

First of ships: \emph{naviguer sur la mer}.

\begin{cousinweb}
\emph{naviguer} is a learned borrowing from Latin \emph{nāvigāre} \textquotedblleft{}to sail\textquotedblright{} (EtymDB; Lewis \& Short), built on \emph{nāvis} \textquotedblleft{}ship\textquotedblright{}.
\end{cousinweb}

\begin{grammarlens}[title={What you've built}]
One more word for telling, explaining and describing at B1.
\end{grammarlens}

\subsection*{Your turn}
\begin{itemize}
\raggedright
  \item \emph{Say it:} \emph{naviguer}
  \item \emph{Say it:} \emph{naviguer}, once more
  \item \emph{Say it:} a sentence of your own with \emph{naviguer}
  \item \emph{Recall:} say the French for the network, then the French for the connection, then say \emph{naviguer} again
\end{itemize}

\begin{tcolorbox}[breakable,colback=teal!4,colframe=teal!35!black,title={Before you move on}]
What does naviguer mean? (\textquotedblleft{}To browse; to sail\textquotedblright{}.) Say it once more.
\end{tcolorbox}
\section[le lien]{le lien — the link; the bond}

\label{lesson:FR-C398-lien}



\begin{quote}
Before the new one: say the French for the connection, then the French for to browse.
\end{quote}

\subsection*{You'll want to know: le lien}
\textbf{le lien} — \textquotedblleft{}the link; the bond\textquotedblright{}.

\begin{quote}
Je t'envoie le lien vers la vidéo.
\end{quote}

\begin{quote}
\emph{I am sending you the link to the video.}
\end{quote}

Between people too: \emph{un lien d'amitié}.

\begin{cousinweb}
\emph{lien} is inherited from Latin \emph{ligāmen} \textquotedblleft{}a bond, a band\textquotedblright{} (EtymDB), from \emph{ligāre} \textquotedblleft{}to tie\textquotedblright{} (Lewis \& Short).
\end{cousinweb}

\begin{grammarlens}[title={What you've built}]
One more word for telling, explaining and describing at B1.
\end{grammarlens}

\subsection*{Your turn}
\begin{itemize}
\raggedright
  \item \emph{Say it:} \emph{le lien}
  \item \emph{Say it:} \emph{le lien}, once more
  \item \emph{Say it:} a sentence of your own with \emph{le lien}
  \item \emph{Recall:} say the French for the connection, then the French for to browse, then say \emph{le lien} again
\end{itemize}

\begin{tcolorbox}[breakable,colback=teal!4,colframe=teal!35!black,title={Before you move on}]
What does le lien mean? (\textquotedblleft{}The link; the bond\textquotedblright{}.) Say it once more.
\end{tcolorbox}
\section[publier]{publier — to post, to publish}

\label{lesson:FR-C398-publier}



\begin{quote}
Before the new one: say the French for to browse, then the French for the link.
\end{quote}

\subsection*{You'll want to know: publier}
\textbf{publier} — \textquotedblleft{}to post, to publish\textquotedblright{}.

\begin{quote}
Elle a publié une photo de son chat.
\end{quote}

\begin{quote}
\emph{She posted a photo of her cat.}
\end{quote}

Of books too: \emph{publier un roman}.

\begin{cousinweb}
\emph{publier} is inherited from Latin \emph{pūblicāre} \textquotedblleft{}to make public\textquotedblright{} (EtymDB; Lewis \& Short), from \emph{pūblicus} \textquotedblleft{}of the people\textquotedblright{}.
\end{cousinweb}

\begin{grammarlens}[title={What you've built}]
One more word for telling, explaining and describing at B1.
\end{grammarlens}

\subsection*{Your turn}
\begin{itemize}
\raggedright
  \item \emph{Say it:} \emph{publier}
  \item \emph{Say it:} \emph{publier}, once more
  \item \emph{Say it:} a sentence of your own with \emph{publier}
  \item \emph{Recall:} say the French for to browse, then the French for the link, then say \emph{publier} again
\end{itemize}

\begin{tcolorbox}[breakable,colback=teal!4,colframe=teal!35!black,title={Before you move on}]
What does publier mean? (\textquotedblleft{}To post, to publish\textquotedblright{}.) Say it once more.
\end{tcolorbox}
